% ============================================================
% 正文：Experimental Setup
% 作者：Kehui Chen, Jacky Keung, Xiaoxue Ma
% 主题：MCP 故障定位（fault localization）根因实证研究
% ============================================================
\section{Experimental Setup}\label{sec:setup}

This section describes the second phase of our sequential-exploratory study. Driven by the
fault-diagnosis pain points identified in the interviews, we build a \textbf{real-bug dataset}
by mining the entire issue tracker of the official MCP SDK, reproduce a representative subset
at runtime over the production Streamable-HTTP transport, and run controlled observability
experiments. We first describe the dataset construction and sampling rationale
(Section~\ref{sec:sampling}), then the four-layer taxonomy (Section~\ref{sec:taxonomy}) and the
reproduction / experiment configuration.

\subsection{Reproduction Setup and Observability Experiments}

We reproduce each of the \textbf{49 active real bugs} over the production
\textbf{Streamable-HTTP + uvicorn} transport, so every case exercises the same transport a real
MCP deployment uses rather than an in-memory substitute. For each bug we write a small trigger
script that builds the relevant MCP server (a subset of the official \texttt{mcp} 2.2.0 SDK),
connects a client, invokes the faulting tool call, and captures the raw SDK return --- either
an error object with \texttt{is\_error} set, a successful result, or a genuine Python / HTTP
exception. Each case is bound to a real issue number (\texttt{issue\_src}), so no fault is
fabricated. This raw return is the evidence base from which we later derive the exposure channel
(Section~\ref{sec:results}): we read the error \emph{text} the SDK actually produced, rather than
relying on any pre-assigned label.

For each execution we record the same four kinds of evidence used throughout the study:
\begin{enumerate}
  \item The raw error or success output returned by the MCP SDK.
  \item The apparent layer at which the SDK exposes the error message.
  \item The real root-cause layer according to our four-layer taxonomy (Section~\ref{sec:taxonomy}).
  \item Whether a silent failure occurs (logical error while \texttt{is\_error=False}).
\end{enumerate}

Two core evaluation metrics are collected:
\begin{itemize}
  \item \textbf{Manifestation-root-cause mismatch}: whether the reported error layer differs
        from the actual fault root-cause layer.
  \item \textbf{Silent failure}: whether erroneous inputs produce an SDK success status
        without explicit error flags.
\end{itemize}

\paragraph{Diagnosis baseline and intervention.} To connect these two metrics to the
practitioners' complaint, we additionally implement a surface-text LLM diagnostic agent
(DeepSeek via the CityU proxy, temperature 0) that reads only the raw error or success text
returned by the SDK for each of the 49 reproduced bugs and predicts the root-cause layer (RQ3).
We first measure its accuracy on the 49 bugs (baseline), then release more control by letting it
rank the four layers (Top-K). We then apply a \emph{runnable diagnostic-clue tool} that reads the
same SDK text, classifies the exposure channel, and emits directional guidance (where to look,
what to verify) without naming the root-cause layer, and re-measure the agent's accuracy to
quantify the value of automated diagnostic guidance over raw surface text.

\paragraph{Experimental limitations.} All experiments were conducted against a single pinned
Python MCP SDK (\texttt{mcp} 2.2.0) over the Streamable-HTTP transport; behaviour may differ
across other language implementations or transports. Although the dataset is a complete census,
the 49 reproduced bugs are a representative subset of the 342 active defects (the remainder
require authentication servers, security-sensitive setups, or other transports that are outside
the scope of a controlled study, as discussed in Section~\ref{sec:sampling}), so our controlled
observations are drawn from that reproducible subset rather than the whole pool.

\subsection{Dataset Construction and Sampling Rationale}
\label{sec:sampling}

Reviewers of fault studies rightly ask three questions of any benchmark: \emph{how large is it},
\emph{how were the faults collected}, and \emph{how can we be sure it represents the cases that
actually occur}. Rather than answer these with an assertion, we describe the data-construction
pipeline we followed and give the resulting counts, so that the provenance and the selection
criteria are explicit and auditable.

\paragraph{Dataset size.} Our dataset is a \emph{complete census} of the publicly reported
defects in one representative, widely deployed MCP implementation: the official
\verb|modelcontextprotocol/python-sdk| (the reference Python SDK, $\approx$24.5k GitHub stars).
We retrieved \emph{every} genuine issue in its tracker (state = all), deliberately excluding
pull requests and non-bug items, giving \textbf{1,343} pure issues spanning the project's full
lifetime (2024--2026). We then manually reviewed every candidate against an explicit set of
exclusion criteria; items were removed only when they were a usage question, a
pure feature/enhancement request, a defect of a \emph{different} SDK (e.g.\ the TypeScript/Node
client), an environment/package/dependency setup problem, a CI/test/refactor item, a duplicate,
or an unverifiable report. This left \textbf{398} confirmed real defects. Because the study pins
a single SDK version (\texttt{mcp} 2.2.0), we then applied a \emph{version-aware} filter:
we reproduced each defect on that pinned version over the real Streamable-HTTP transport and
removed the \textbf{56} that are already \emph{fixed in 2.2.0} (i.e.\ the reported misbehaviour
no longer occurs, so they are not live defects of the studied version). This leaves an active
dataset of \textbf{342} defects that genuinely manifest on the studied SDK. Because this is a
census rather than a sample, the dataset is \emph{complete with respect to the library}: there
is no selection gap between the repository's issue history and the faults we study.

% ★【中文说明】新增：数据集筛选漏斗表。从官方 python-sdk 全量 issue(1343) 出发，人工复核确认真 bug(398)，
%   减去 2.2.0 已修复的(56)，得有效数据集(342)，其中 49 条可复现、293 条不可复现。
% ★[English] New: dataset-screening funnel table. From the full issue pool (1,343) through
%   manual confirmation of real bugs (398), minus the 56 already fixed in 2.2.0, giving the
%   342 active dataset, of which 49 are reproducible and 293 are not.
Table~\ref{tab:funnel} summarises this screening funnel.
\begin{table}[!t]
\centering
\caption{Data-collection funnel for the real-bug dataset mined from the official \texttt{modelcontextprotocol/python-sdk} issue tracker.}
\label{tab:funnel}
\scriptsize
\renewcommand{\arraystretch}{1.05}
\setlength{\tabcolsep}{2.5pt}
\begin{tabularx}{\columnwidth}{@{}>{\centering\arraybackslash}l >{\centering\arraybackslash}c >{\centering\arraybackslash}X@{}}
\toprule
\textbf{Stage} & \textbf{\#} & \textbf{Step} \\
\midrule
1 & 1{,}343 & Official pure issues in the tracker (pull requests excluded) \\
2 & 398 & Confirmed real bugs (manual review) \\
3 & $-$56 & Already fixed in the pinned SDK (\texttt{mcp} 2.2.0), removed \\
\midrule
\textbf{4} & \textbf{342} & \textbf{ACTIVE dataset (studied)} \\
\midrule
5 & \textbf{49} & \textbf{Reproduced at runtime (evidence-based)} \\
6 & 293 & Not reproducible at runtime (see Table~\ref{tab:notrepro}) \\
\bottomrule
\end{tabularx}
\vspace{0.3em}
\parbox{\columnwidth}{\scriptsize\centering
\textit{Note:} Stage 4 subtracts the defects that behave correctly on the pinned SDK version
(\texttt{mcp} 2.2.0). Stage 6 bugs are each bound to a real issue number and reproduced over the
production Streamable-HTTP transport.}
\end{table}

\paragraph{Collection criteria.} Three criteria governed the collection. First,
\emph{completeness}: we retrieved all issues of the library, not a convenience subset, so no
fault category is silently absent. Second, \emph{scoping}: only genuine defects of the Python
SDK were retained; questions, feature requests, other-SDK faults, and environmental/packaging
issues were excluded, so the dataset isolates SDK-internal behaviour. Third,
\emph{reproducibility}: we pin a single SDK version (\texttt{mcp} 2.2.0) and reproduce a
representative subset of the defects over the real Streamable-HTTP transport, so that the
faults are not merely textual reports but executable, end-to-end cases whose error-reporting
behaviour can be observed and measured.

\paragraph{How representativeness is guaranteed.} Because we take a complete census, the
dataset does not suffer the sample-selection bias of a partial scrape. Moreover, the 342
active defects span the whole system surface rather than a single component: applying our
four-layer origin taxonomy, they break down as \textbf{Server (91)}, \textbf{Client (118)},
\textbf{Transport (85)}, and \textbf{Host (48)}, and cover both open and closed reports (93
open / 249 closed). This breadth is consistent with the four-layer taxonomy we use for analysis,
ensuring the dataset does not over-represent a single layer. Finally, a reproduced subset is
sampled against the true frequency of the reported defects, so the controllable experiments
inherit the distribution of the real world rather than an arbitrary one.

\paragraph{Why a single library is sufficient.} Taraghi et al.~\cite{taraghi2026real} mined a
large \emph{cross-library} corpus (13,555 repositories, 3,282 bug issues), and so report more
defects in aggregate than we do. We do not compete on raw volume; the two corpora answer
different questions. Taraghi's breadth establishes the \emph{landscape} of MCP fault types
across many heterogeneous implementations. Our study instead needs \emph{depth and control}: we
measure \emph{how the SDK reports (or hides)} a defect, which is implementation-specific
error-reporting behaviour that requires a pinned, reproducible setup with a real transport.
For this, cross-library breadth is not only unnecessary but harmful: merging many SDK versions
and transports would confound the error-reporting signal with implementation differences.

A single canonical library suffices for this purpose for three reasons. First, the official
Python SDK is the reference implementation and is among the most widely adopted, so it is the
most representative \emph{single} choice of where real MCP faults manifest. Second, our
analysis dimension (the four-layer origin taxonomy and the three exposure channels) is
\emph{SDK-agnostic} and does not rely on library-specific features; it can be applied to any
implementation, so the mechanisms we quantify are not artefacts of one library. Third,
controlled reproducibility (a pinned version plus the real transport) is only feasible within
a coherent, single-version codebase; a cross-library corpus would make the controlled
observability experiments impossible to run consistently. We therefore deliberately trade
breadth for depth, and we make this trade-off explicit so that readers can weigh it: the
dataset is a complete census of one canonical library rather than a partial sample of many.

\paragraph{Why the 293 bugs were not reproduced.} Of the 342 active real bugs, we reproduced
49 at runtime and formed the evidence-based analysis. The remaining \textbf{293} could not be
reproduced over the production Streamable-HTTP transport; the reason for each was recorded
during reproduction and is grouped into paper-level categories in Table~\ref{tab:notrepro}. The
dominant causes are environment / transport configuration (106, 36\%), authentication
credentials (48, 16\%), and defects that a single clean tool call does not exercise (88, 30\%,
e.g.\ concurrency races, cross-server state, or client-side observation); a smaller set is
purely observational or scenario-specific (41, 14\%) or reflects a protocol specification gap
(6) or an external-tool / inspector dependency (4). We report each of these honestly rather than
asserting the bug is fixed, and we never fabricate an exposure channel for a bug we could not
reproduce.

% ★【中文说明】新增：293 条无法复现 bug 的原因分类表（环境/认证/未触发/观察型/规范缺口/外部工具）。
% ★[English] New: category table for the 293 not-reproducible bugs (env/transport, auth,
%   not-exercised, observation, spec-gap, external-tool).
\begin{table}[!t]
\centering
\caption{Why the 293 active real bugs were not reproduced at runtime.}
\label{tab:notrepro}
\scriptsize
\renewcommand{\arraystretch}{1.05}
\setlength{\tabcolsep}{2.5pt}
\begin{tabularx}{\columnwidth}{@{}>{\raggedright\arraybackslash}X >{\centering\arraybackslash}c >{\centering\arraybackslash}c@{}}
\toprule
\textbf{Reason (not reproducible)} & \textbf{\#} & \textbf{\% of 293} \\
\midrule
Environment / transport configuration & 106 & 36.2\% \\
Authentication / credentials required & 48 & 16.4\% \\
Not exercised / no defect captured by a single call & 88 & 30.0\% \\
Observation / scenario-specific & 41 & 14.0\% \\
Protocol specification gap & 6 & 2.0\% \\
External tool / inspector dependency & 4 & 1.4\% \\
\midrule
\textbf{Total} & \textbf{293} & \textbf{100\%} \\
\bottomrule
\end{tabularx}
\end{table}

\paragraph{Honesty of the provenance.} We do not claim to reproduce every one of the 342
active defects: reproducing all of them would require authentication servers, security-sensitive
setups, and other transports/SDK versions outside the scope of a controlled study. We reproduce
a representative subset (49 distinct, verified over the real Streamable-HTTP transport), and we
report the remaining defects as confirmed-but-not-reproduced so that the distinction between
"confirmed defect" and "executable case" is transparent. Where we map an issue to a Taraghi
category, the label is ours (we place each report into Taraghi's taxonomy buckets); it is
\emph{not} the raw per-label count of Taraghi et al.~\cite{taraghi2026real}, and we say so
explicitly rather than presenting an author-labelled mapping as their statistics.